\documentclass[14pt,a4paper]{extarticle}
\usepackage{fontspec}
\setmainfont{Times New Roman}
\usepackage[russian]{babel}
\usepackage[left=30mm,right=15mm,top=20mm,bottom=20mm]{geometry}
\usepackage{setspace}
\onehalfspacing
\renewcommand{\normalsize}{\fontsize{14pt}{16.8pt}\selectfont}
\normalsize
\usepackage{indentfirst}
\setlength{\parindent}{1.25cm}
\setlength{\parskip}{0pt}
\usepackage{enumitem}
\setlist{nosep,leftmargin=1.25cm}
\usepackage{titlesec}
\titleformat{\section}{\normalfont\fontsize{14}{17}\selectfont\bfseries}{\thesection}{1em}{}
\titleformat{\subsection}{\normalfont\fontsize{14}{17}\selectfont\bfseries}{\thesubsection}{1em}{}
\titlespacing*{\section}{0pt}{14pt}{14pt}
\titlespacing*{\subsection}{0pt}{12pt}{8pt}
\usepackage{tocloft}
\renewcommand{\contentsname}{Содержание}
\renewcommand{\cfttoctitlefont}{\normalfont\fontsize{14}{17}\selectfont\bfseries}
\renewcommand{\cftsecfont}{\normalfont}
\renewcommand{\cftsecpagefont}{\normalfont}
\setlength{\cftbeforesecskip}{0pt}
\usepackage{longtable,array}
\usepackage{tikz}
\usetikzlibrary{arrows.meta,positioning,shapes.multipart,calc}
\usepackage{caption}
\DeclareCaptionLabelSeparator{hyphen}{ - }
\captionsetup{font=normalsize,labelsep=hyphen}
\addto\captionsrussian{\renewcommand{\figurename}{Рисунок}}
\usepackage[hidelinks]{hyperref}
\hypersetup{pdftitle={Лабораторная работа №1. Вариант 84121905},pdfauthor={Лейковский Никита}}
\pagestyle{plain}
\emergencystretch=2em
\tikzset{cls/.style={rectangle split,rectangle split parts=2,draw,align=left,font=\fontsize{11}{13}\selectfont,inner sep=6pt}, dep/.style={dashed,-{Stealth}}, assoc/.style={-{Stealth}}, pkg/.style={draw,minimum width=42mm,minimum height=14mm,font=\fontsize{12}{14}\selectfont}}
\begin{document}
\begin{titlepage}
\begin{center}
Университет ИТМО\\
Информационные системы
\vfill
\textbf{Отчёт по лабораторной работе №1}\\[10pt]
Управление транспортными средствами\\
Вариант 84121905
\end{center}
\vfill
\begin{flushright}
Выполнил: Лейковский Никита\\
ИСУ 466497
\end{flushright}
\vfill
\begin{center}Санкт-Петербург\\2026\end{center}
\end{titlepage}
\setcounter{page}{2}
\tableofcontents
\clearpage
\section{Текст задания}
Реализовать информационную систему, которая позволяет взаимодействовать с объектами класса Vehicle, описание которого приведено ниже.

\subsection*{Модель данных}
Класс Vehicle содержит следующие закрытые поля. Типы и ограничения соответствуют варианту 84121905.

{\setstretch{1.0}
\begin{longtable}{|p{.25\textwidth}|p{.20\textwidth}|p{.45\textwidth}|}
\hline
Поле & Тип & Ограничения\\\hline
\endhead
id & Integer & Не null; больше 0; уникальное; генерируется автоматически.\\\hline
name & String & Не null; строка не может быть пустой.\\\hline
coordinates & Coordinates & Не null.\\\hline
creationDate & java.util.Date & Не null; генерируется автоматически.\\\hline
type & VehicleType & Не null.\\\hline
enginePower & int & Больше 0.\\\hline
numberOfWheels & Integer & Не null; больше 0.\\\hline
capacity & Double & Не null; больше 0.\\\hline
distanceTravelled & int & Больше 0.\\\hline
fuelConsumption & double & Больше 0.\\\hline
fuelType & FuelType & Не null.\\\hline
\end{longtable}}
Класс Coordinates содержит закрытые поля x типа Double и y типа Float. Оба поля не могут быть null. Максимальное значение y: 408.

Перечисление VehicleType: PLANE, BOAT, SHIP. Перечисление FuelType: DIESEL, ALCOHOL, MANPOWER.

\subsection*{Основные требования}
\begin{enumerate}
\item Основное назначение информационной системы - управление объектами, созданными на основе заданного в варианте класса.
\item Необходимо, чтобы с помощью системы можно было выполнить следующие операции с объектами: создание нового объекта, получение информации об объекте по ИД, обновление объекта (модификация его атрибутов), удаление объекта. Операции должны осуществляться в отдельных окнах (интерфейсах) приложения. При получении информации об объекте класса должна также выводиться информация о связанных с ним объектах.
\item При создании объекта класса необходимо дать пользователю возможность связать новый объект с объектами вспомогательных классов, которые могут быть связаны с созданным объектом и уже есть в системе.
\item Выполнение операций по управлению объектами должно осуществляться на серверной части (не на клиенте), изменения должны синхронизироваться с базой данных.
\item На главном экране системы должен выводиться список текущих объектов в виде таблицы (каждый атрибут объекта - отдельная колонка в таблице). При отображении таблицы должна использоваться пагинация (если все объекты не помещаются на одном экране).
\item Нужно обеспечить возможность фильтровать/сортировать строки таблицы, которые показывают объекты (по значениям любой из строковых колонок). Фильтрация элементов должна производиться только по полному совпадению.
\item Переход к обновлению (модификации) объекта должен быть возможен из таблицы с общим списком объектов и из области с визуализацией объекта (при её реализации).
\item При добавлении/удалении/изменении объекта он должен автоматически появиться/исчезнуть/измениться в интерфейсах у других пользователей, авторизованных в системе.
\item Если при удалении объекта с ним связан другой объект, связанные объекты должны быть связаны с другим объектом (по выбору пользователя), а изначальный объект удалён.
\item Пользователи должны иметь возможность просмотра всех объектов. Для модификации объекта должно открываться отдельное диалоговое окно. При вводе некорректных значений в поля объекта должны появляться информативные сообщения о соответствующих ошибках.
\end{enumerate}

\subsection*{Специальные операции}
В системе должен быть реализован отдельный пользовательский интерфейс для выполнения специальных операций над объектами:
\begin{enumerate}
\item Удалить все объекты, значение поля fuelType которого эквивалентно заданному.
\item Удалить один (любой) объект, значение поля fuelConsumption которого эквивалентно заданному.
\item Вернуть количество объектов, значение поля fuelConsumption которых больше заданного.
\item Найти все транспортные средства заданного типа.
\item Найти все транспортные средства с мощностью двигателя в заданном диапазоне.
\end{enumerate}
Представленные операции должны быть реализованы в рамках компонентов бизнес-логики приложения без прямого использования функций и процедур БД.

\subsection*{Хранение объектов}
\begin{enumerate}
\item Организовать хранение данных об объектах в реляционной СУБД PostgreSQL. Каждый объект, с которым работает ИС, должен быть сохранён в базе данных.
\item Все требования к полям класса, указанные в виде комментариев к описанию классов, должны быть выполнены на уровне ORM и БД.
\item Для генерации поля id использовать средства базы данных.
\item Для подключения к БД на кафедральном сервере использовать хост pg, имя базы данных - studs, имя пользователя/пароль совпадают с таковыми для подключения к серверу.
\end{enumerate}

\subsection*{Взаимодействие с пользователем}
\begin{enumerate}
\item Система должна реагировать на некорректный пользовательский ввод, ограничивая ввод недопустимых значений и информируя пользователей о причине ошибки.
\item Переходы между различными логически обособленными частями системы должны осуществляться с помощью меню.
\item При добавлении/удалении/изменении объекта он должен автоматически появиться/исчезнуть/измениться на области у всех других клиентов.
\end{enumerate}

\subsection*{Требования к разработке}
\begin{enumerate}
\item В качестве основы для реализации ИС необходимо использовать Spring MVC.
\item Для создания уровня хранения необходимо использовать Hibernate.
\item Разные уровни приложения должны быть отделены друг от друга, разные логические части ИС должны находиться в отдельных компонентах.
\end{enumerate}

\subsection*{Состав отчёта}
Текст задания; UML-диаграммы классов и пакетов разработанного приложения; исходный код системы или ссылка на репозиторий с исходным кодом; выводы по работе. Репозиторий при сдаче должен быть публичным; если ветка не указана, используется основная.

\clearpage
\section{UML-диаграммы}
\subsection*{Классы предметной области}
\begin{center}
\begin{tikzpicture}
\node[cls,text width=68mm] (v) {\centering\textbf{Vehicle}\par $\langle\!\langle$Entity$\rangle\!\rangle$\nodepart{second}
- id: Integer\\- name: String\\- coordinates: Coordinates\\- creationDate: Date\\- type: VehicleType\\- enginePower: int\\- numberOfWheels: Integer\\- capacity: Double\\- distanceTravelled: int\\- fuelConsumption: double\\- fuelType: FuelType\\[4pt]+ update(...): void};
\node[cls,text width=40mm,right=18mm of v.north east,anchor=north west] (c) {\centering\textbf{Coordinates}\par $\langle\!\langle$Entity$\rangle\!\rangle$\nodepart{second}- id: Integer\\- x: Double\\- y: Float\\[4pt]+ update(x, y): void};
\node[cls,text width=40mm,below=22mm of c] (t) {\centering\textbf{VehicleType}\par $\langle\!\langle$enumeration$\rangle\!\rangle$\nodepart{second}PLANE\\BOAT\\SHIP};
\node[cls,text width=40mm,below=16mm of t] (f) {\centering\textbf{FuelType}\par $\langle\!\langle$enumeration$\rangle\!\rangle$\nodepart{second}DIESEL\\ALCOHOL\\MANPOWER};
\draw[assoc] (v.east |- c.west) -- node[above,font=\small]{0..* \hspace{6mm} 1} (c.west);
\draw[assoc] (v.east |- t.west) -- (t.west);
\draw[assoc] (v.south) -- ++(0,-8mm) |- (f.west);
\end{tikzpicture}
\captionof{figure}{Классы предметной области}
\end{center}
Coordinates является отдельной сущностью. Несколько объектов Vehicle могут ссылаться на один набор координат через внешний ключ coordinates\_id. Стрелки обозначают направленные ассоциации. Методы доступа на диаграмме опущены.

\clearpage
\subsection*{Классы приложения}
\begin{center}
\begin{tikzpicture}[node distance=16mm and 16mm]
\node[cls,text width=61mm] (vc) {\centering\textbf{VehicleController}\nodepart{second}+ list(...): String\\+ create(...): String\\+ update(...): String\\+ delete(...): String};
\node[cls,text width=61mm,right=of vc] (sc) {\centering\textbf{SpecialOperationsController}\nodepart{second}+ deleteByFuel(...): String\\+ deleteOneByConsumption(...): String\\+ count(...): String\\+ byType(...): String\\+ byPower(...): String};
\node[cls,text width=61mm,below=of vc] (vs) {\centering\textbf{VehicleService}\nodepart{second}+ findPage(...): Page\\+ get(id): Vehicle\\+ create(form): Vehicle\\+ update(id, form): Vehicle\\+ delete(id): void};
\node[cls,text width=61mm,below=of sc] (ss) {\centering\textbf{SpecialOperationsService}\nodepart{second}+ deleteAllByFuelType(...): long\\+ deleteOneByFuelConsumption(...)\\+ countByFuelConsumption\\\hspace*{3mm}GreaterThan(...): long\\+ findByType(...): List\\+ findByEnginePowerRange(...): List};
\node[cls,text width=65mm,below=22mm of $(vs.south)!0.5!(ss.south)$] (r) {\centering $\langle\!\langle$interface$\rangle\!\rangle$\par\textbf{VehicleRepository}\nodepart{second}JpaRepository<Vehicle, Integer>\\JpaSpecificationExecutor<Vehicle>};
\draw[assoc] (vc)--(vs);\draw[assoc] (sc)--(ss);
\draw[assoc] (vs.south) |- ($(r.north)+(0,7mm)$) -- (r.north);
\draw[assoc] (ss.south) |- ($(r.north)+(0,7mm)$) -- (r.north);
\end{tikzpicture}
\captionof{figure}{Контроллеры, сервисы и репозиторий}
\end{center}
Контроллер VehicleController получает данные через DTO VehicleForm. Сервис VehicleService использует этот DTO при создании и изменении объектов. Параметры методов сокращены для читаемости. VehicleRepository расширяет интерфейсы Spring Data JPA.

\clearpage
\subsection*{Связанные объекты и вход}
\begin{center}
\begin{tikzpicture}[node distance=16mm and 14mm]
\node[cls,text width=62mm] (cc) {\centering\textbf{CoordinatesController}\nodepart{second}+ list(model): String\\+ create(x, y, ...): String\\+ update(id, x, y, ...): String\\+ delete(id, replacementId, ...): String};
\node[cls,text width=62mm,right=of cc] (vf) {\centering\textbf{VehicleForm}\par $\langle\!\langle$DTO$\rangle\!\rangle$\nodepart{second}name, coordinatesId\\coordinateX, coordinateY\\type, enginePower\\numberOfWheels, capacity\\distanceTravelled\\fuelConsumption, fuelType};
\node[cls,text width=62mm,below=of cc] (cs) {\centering\textbf{CoordinatesService}\nodepart{second}+ all(): List\\+ get(id): Coordinates\\+ create(x, y): Coordinates\\+ update(id, x, y): void\\+ delete(id, replacementId): void};
\node[cls,text width=62mm,below=of vf] (cfg) {\centering\textbf{SecurityConfig}\nodepart{second}+ security(http)\\\hspace*{3mm}: SecurityFilterChain\\+ passwords(): PasswordEncoder};
\node[cls,text width=62mm,below=of cs] (cr) {\centering $\langle\!\langle$interface$\rangle\!\rangle$\par\textbf{CoordinatesRepository}\nodepart{second}JpaRepository<Coordinates, Integer>};
\node[cls,text width=62mm,below=of cfg] (lc) {\centering\textbf{LoginController}\nodepart{second}+ login(): String\\+ registration(model): String\\+ register(form, errors): String};
\draw[assoc] (cc)--(cs);\draw[assoc] (cs)--(cr);
\end{tikzpicture}
\captionof{figure}{Координаты, DTO и конфигурация входа}
\end{center}
CoordinatesService также использует VehicleRepository для поиска и переназначения связанных объектов. VehicleService обращается к CoordinatesService при выборе либо создании координат. SecurityConfig задаёт обязательный вход; LoginController обрабатывает вход и регистрацию.

\clearpage
\subsection*{Регистрация и учётные записи}
\begin{center}
\begin{tikzpicture}[node distance=18mm and 14mm]
\node[cls,text width=62mm] (lc) {\centering\textbf{LoginController}\nodepart{second}+ register(form, errors): String};
\node[cls,text width=62mm,right=of lc] (rf) {\centering\textbf{RegistrationForm}\nodepart{second}- username: String\\- password: String\\- confirmation: String\\+ isPasswordsMatch(): boolean};
\node[cls,text width=62mm,below=of lc] (us) {\centering\textbf{UserService}\nodepart{second}UserDetailsService\\+ register(username, password): void\\+ loadUserByUsername(username)\\\hspace*{3mm}: UserDetails\\+ ensureUser(username, password)};
\node[cls,text width=62mm,below=of rf] (ua) {\centering\textbf{UserAccount}\nodepart{second}- id: Long\\- username: String\\- passwordHash: String};
\node[cls,text width=62mm,below=of us] (ur) {\centering\textbf{UserAccountRepository}\nodepart{second}JpaRepository<UserAccount, Long>\\+ findByUsername(username)\\+ existsByUsername(username)};
\draw[assoc] (lc)--(us);
\draw[dep] (lc)--(rf);
\draw[assoc] (us)--(ur);
\draw[dep] (ur.east) -| (ua.south);
\end{tikzpicture}
\captionof{figure}{Регистрация и хранение пользователей}
\end{center}
Учётные записи хранятся в PostgreSQL. Логин уникален и приводится к нижнему регистру. Пароли хешируются BCrypt. StartupData создаёт начальные учётные записи и 24 примера транспорта; SeedState отмечает выполненное заполнение, исключая повторное добавление.

\clearpage
\subsection*{Обновление клиентов}
\begin{center}
\begin{tikzpicture}[node distance=18mm and 14mm]
\node[cls,text width=60mm] (v) {\centering\textbf{VehicleService}\nodepart{second}CRUD-операции};
\node[cls,text width=60mm,right=of v] (s) {\centering\textbf{SpecialOperationsService}\nodepart{second}Специальные операции};
\node[cls,text width=66mm,below=20mm of $(v.south)!0.5!(s.south)$] (e) {\centering\textbf{VehicleChangedEvent}\par $\langle\!\langle$record$\rangle\!\rangle$\nodepart{second}operation: String\\vehicleId: Integer};
\node[cls,text width=66mm,below=of e] (l) {\centering\textbf{VehicleChangeListener}\nodepart{second}+ onVehicleChanged(event): void};
\node[cls,text width=66mm,below=of l] (h) {\centering\textbf{VehicleWebSocketHandler}\nodepart{second}+ broadcast(event): void};
\draw[dep] (v.south) |- ($(e.north)+(0,7mm)$) -- (e.north);
\draw[dep] (s.south) |- ($(e.north)+(0,7mm)$) -- (e.north);
\draw[dep] (l) -- node[right,font=\small]{получает событие} (e);
\draw[assoc] (l) -- (h);
\end{tikzpicture}
\captionof{figure}{Классы синхронизации}
\end{center}
Сервисы, включая CoordinatesService, публикуют событие через механизм событий Spring. Listener обрабатывает его после фиксации транзакции. WebSocketConfig регистрирует обработчик по адресу /ws/vehicles. Пунктир обозначает зависимость.

\clearpage
\subsection*{Пакеты приложения}
\begin{center}
\begin{tikzpicture}[node distance=20mm and 25mm]
\node[pkg] (web) {web};
\node[pkg,right=of web] (srv) {service};
\node[pkg,below=of srv] (repo) {repository};
\node[pkg,below=of web] (rt) {realtime};
\node[pkg,below=of repo] (dom) {domain};
\node[pkg,below=of rt] (sec) {security};
\draw[dep] (web) -- (srv);
\draw[dep] (srv) -- (repo);
\draw[dep] (repo) -- (dom);
\draw[dep] (web.east) -- ++(9mm,0) |- (dom.west);
\draw[dep] (srv.east) -- ++(8mm,0) |- (dom.east);
\draw[dep] (rt) -- (srv);
\draw[dep] (srv.north) -- ++(0,10mm) -| node[pos=.75,above,font=\small]{DTO} (web.north);
\end{tikzpicture}
\captionof{figure}{Зависимости пакетов ru.itmo.vehiclelab}
\end{center}
Пакет web содержит MVC-контроллеры, VehicleForm и обработчик ошибок; service - бизнес-логику, событие и исключение; repository - доступ к данным; domain - модель; realtime - WebSocket-компоненты; security - конфигурацию входа и защиты запросов.

Зависимость service от web обусловлена использованием VehicleForm. Пакет realtime использует VehicleChangedEvent из service. Hibernate выполняет отображение модели в PostgreSQL; HTML-шаблоны и статические ресурсы находятся отдельно от Java-пакетов.

\clearpage
\section{Исходный код}
Исходный код и конфигурация запуска опубликованы в репозитории:
\begin{center}
\url{https://github.com/wonkersone/informatsionnye-sistemy-lr1}
\end{center}

Приложение использует Java 21, Spring MVC, Hibernate, Spring Data JPA, Thymeleaf и Flyway. Основной профиль подключается к PostgreSQL. Профиль demo использует временную H2-базу и предназначен только для демонстрации. Адрес локального интерфейса: \url{http://localhost:6767}.

Вход реализован средствами Spring Security; формы защищены CSRF-токенами. Координаты выбираются из существующих наборов либо создаются вместе с транспортом. При удалении используемых координат связанные транспортные средства переназначаются на выбранный набор в одной транзакции. Редактирование транспорта открывается в отдельном диалоговом окне.

\section{Выводы}
Разработана информационная система на Spring MVC и Hibernate с CRUD, пятью специальными операциями, регистрацией, авторизацией и синхронизацией клиентов. Реализованы хранение в PostgreSQL, проверка ограничений и переназначение связей.
\end{document}
